\PageHeading{13}{已知限制与迭代计划}{优先补齐影响正式参赛判定的证据，再扩展部署能力}
\Sub{当前边界}
\begin{itemize}
\item \Lead{知识来源验收。}ima读取额度限制导致新流程完整验收尚未完成。本文本地样例不计作ima来源证明；需在额度恢复和权限确认后重跑正式任务。
\item \Lead{语言与视觉质量。}规则演示仍有清单化、参数化及跨渠道句式复用；多个主题主要改变颜色/比例，尚非成熟的独立构图体系。模型也会留下待确认问题，当前缺少真实客户样本盲评。
\item \Lead{运行与权限。}当前是单操作者本机Demo，无多用户鉴权、持久数据库及公开服务的部署配置。任务/大纲在内存中，重启后失效。
\item \Lead{发布与素材。}公众号图尚未转存微信；没有账号自动发布、印刷PDF与在线图层编辑。实际微信粘贴和平台发布效果未验收。
\item \Lead{资料格式。}当前适配原文以Markdown、文本、HTML及有权限笔记为主；PDF/Office原文在连接器中明确提示不支持。
\end{itemize}
\Sub{按验收价值排序的工作项}
\begingroup\small
\begin{tabularx}{\linewidth}{P{12mm}Y Y}
\rowcolor{Forest}\color{white}\bfseries 优先级 & \color{white}\bfseries 下一步工作 & \color{white}\bfseries 完成判据\\
P0 & 恢复并核验ima权限/额度，按两类场景、两平台完成至少3组正式样例。 & 每组保存事实与原图来源、真实调用记录、核验结论和可访问成品。\\
\rowcolor{Mist}P0 & 补齐魔搭页面及真实运行录屏，统一各提交入口的信息。 & 评审期无需额外授权可访问；视频连续展示任务到导出，并补换场景/平台。\\
P0 & 对统一任务做人工核验与对抗测试，处理标题等非数值断言。 & 混淆型号、无依据价格/补贴、伪造体验、缺字段及临时修改有明确结果。\\
\rowcolor{Mist}P1 & 建立可重复评估集，比较模型输出与人工销售文案。 & 给出样本量、评价口径、质量与耗时分布；不以测试通过率代替准确率。\\
P1 & 进一步优化渠道语言、提示词与长条件排版。 & 人工评分和溢出统计改善，限定条件保留且可读。\\
\rowcolor{Mist}P2 & 多用户权限、持久化、队列与部署加固；按需接平台素材能力。 & 隔离测试、恢复测试、负载测试及安全验收通过后再公开部署。\\
\bottomrule
\end{tabularx}
\endgroup
\begin{note}\small
\Lead{建议评估指标}\quad 无依据断言数、必要条件保留率、图像型号一致率、人工可发布评分、每任务人工修改次数、端到端耗时。均需配套抽样与标注方法；本说明书不为尚未测量的指标填入推测值。
\end{note}
\SourceNote{赛题手册第8--9页与当前代码/实测限制。P0/P1/P2为项目优先级，非赛事评分等级。}
